% ============================================================
% section7_optimization.tex - Раздел 7: Угловая зависимость и интерактивное сравнение с теорией
% ============================================================

\section{Сравнение теории с экспериментом и интерактивное моделирование}

В предыдущих разделах мы подготовили базу: научились рассчитывать экспериментальные спектры пропускания для различных углов и реализовали электродинамическую модель Бланко в модуле \texttt{theoretical.py}. Теперь наша главная цель — сопоставить теорию с экспериментом на конкретной частоте.

\subsection{Вектор Джонса и физические величины}
Напомним, что в нашем эксперименте исходный терагерцовый пучок линейно поляризован. Для теоретических расчетов с использованием матриц Джонса мы задаем начальный вектор Джонса для линейной горизонтальной поляризации:
\begin{equation}
\mathbf{E}_{in} = \begin{pmatrix} 1 \\ 0 \end{pmatrix}
\end{equation}

После прохождения системы поляризаторов мы получаем выходное поле $\mathbf{E}_{out}$. Важно понимать, что угловую зависимость можно анализировать двумя способами:
\begin{itemize}
    \item \textbf{Амплитудное пропускание} ($T_{amp} = |\mathbf{E}_{out}| / |\mathbf{E}_{in}|$): напрямую связано с электрическим полем, измеряемым в THz-TDS.
    \item \textbf{Энергетическое пропускание (интенсивность)} ($T_{pow} = |\mathbf{E}_{out}|^2 / |\mathbf{E}_{in}|^2$): чаще используется для оценки эффективности оптических элементов.
\end{itemize}
В нашем коде мы будем оперировать энергетическим пропусканием (по мощности), так как именно оно характеризует реальное ослабление пучка.

\subsection{Два взгляда на данные: линейная и логарифмическая (dB) шкалы}

Сравнивать теорию с экспериментом мы будем сразу на двух графиках (в виде \texttt{plot grid}):
\begin{enumerate}
    \item \textbf{В процентах (линейная шкала):} Этот график идеально показывает согласование при \textbf{высоком пропускании} (когда поляризаторы параллельны, $\theta = 0^\circ, 180^\circ$). Однако в области скрещенных поляризаторов ($\theta \approx 90^\circ$) значения близки к нулю, и детали теряются.
    \item \textbf{В децибелах (логарифмическая шкала):} Переход к децибелам осуществляется по формуле $T_{dB} = 10 \log_{10}(T_{pow})$. Децибелы растягивают область малых значений, позволяя детально рассмотреть согласование при \textbf{низком пропускании} (в зоне максимальной экстинкции). Падение на $-30$ дБ означает ослабление мощности в 1000 раз!
\end{enumerate}

\subsection{Интерактивная подгонка: зачем нужны слайдеры}

Для расчета теоретической кривой нам нужны параметры решетки. Мы могли бы добавить их в наш файл \texttt{config.py} (например, \texttt{P = 12.5e-6}, \texttt{D = 5.0e-6}, \texttt{PERIOD\_SCALE = 1.0}, \texttt{LOSS\_FACTOR = 0.0}), считывать их оттуда, строить график и каждый раз при несовпадении менять значения в конфигурационном файле, перезапуская скрипт.

Но этот процесс крайне неудобен и долог! Намного эффективнее один раз написать интерактивный интерфейс. Мы воспользуемся виджетом \texttt{Slider} из \texttt{matplotlib.widgets}, чтобы в реальном времени вращать ручки параметров и наблюдать, как теоретическая кривая «натягивается» на экспериментальные точки.

\subsection{Извлечение экспериментальной угловой зависимости}

Чтобы наложить экспериментальные результаты на теоретические кривые, нам нужно извлечь зависимость пропускания от угла поворота для \textit{одной конкретной частоты} (например, 0.5 ТГц). Это делается очень изящно:
\begin{enumerate}
    \item Наш конвейер \texttt{calculate\_transmission} уже на этапе обработки единичного файла делит спектральную амплитуду сигнала $E_{\text{sig}}(\theta)$ на спектральную амплитуду фонового сигнала свободного пространства $E_{\text{bg}}$ и возводит отношение в квадрат.
    \item На выходе мы получаем спектр пропускания по мощности $T_{\text{exp}}(\theta) = \left| E_{\text{sig}}(\theta) / E_{\text{bg}} \right|^2$ относительно открытого терагерцового пучка.
    \item Нам остаётся лишь выбрать значения на интересующей частоте для каждого угла поворота и передать их функции построения графиков.
\end{enumerate}
Эту логику реализует функция \texttt{get\_experimental\_transmission\_at\_freq}, с которой мы познакомимся в листинге ниже.

\subsection{Программная реализация интерактивного окна с оценкой погрешности}

Для наглядного сопоставления теории и эксперимента мы воспользуемся функцией интерактивной подгонки \texttt{run\_step3\_interactive\_fit}, встроенной в основное приложение. Чтобы сделать сравнение строгим и научно обоснованным, мы дополним графики визуализацией остатков и расчётом числовой метрики качества подгонки.

Полный листинг интегрированного приложения \texttt{main.py} вынесен в \textbf{\textit{Приложение~\ref{app:integrated_app}}}. Давай подробно разберём, какие математические инструменты в него заложены.

\subsection{Объективная оценка отклонения: остатки и метрика RMSE}

Просто «на глаз» оценить, насколько хорошо теоретическая кривая совпадает с экспериментальными точками, невозможно — человеческое зрение склонно преувеличивать совпадения в одних областях и игнорировать расхождения в других. Для объективного анализа нам понадобятся две статистические концепции: остатки (residuals) и среднеквадратическая ошибка (RMSE).

\textbf{1. Остатки (Residuals)} характеризуют локальное отклонение модели от эксперимента в каждой измеренной точке $\theta_i$. Мы будем вычислять их отдельно для двух шкал:
\begin{itemize}
    \item \textbf{В процентном представлении (линейный масштаб):}
    \begin{equation}
    R_{\%}(\theta_i) = \left( T_{\text{exp}}(\theta_i) - T_{\text{theo}}(\theta_i) \right) \times 100\%
    \end{equation}
    Этот показатель наглядно демонстрирует абсолютную ошибку в области высокой прозрачности (когда поляризаторы параллельны, $\theta \approx 0^\circ, 90^\circ$).
    
    \item \textbf{В децибелах (логарифмический масштаб):}
    \begin{equation}
    R_{\text{dB}}(\theta_i) = T_{\text{exp, dB}}(\theta_i) - T_{\text{theo, dB}}(\theta_i)
    \end{equation}
    Логарифмические остатки критически важны для оценки качества модели в области глубокой экстинкции (скрещенные поляризаторы, $\theta \approx 45^\circ$). Ошибка всего в $1\%$ в линейной шкале вблизи нуля может превратиться в огромные $10$ дБ расхождения в логарифмической шкале!
\end{itemize}

\textbf{2. Среднеквадратическая ошибка (RMSE, Root Mean Square Error)} объединяет все локальные отклонения в единую интегральную оценку качества нашей модели:
\begin{equation}
\text{RMSE} = \sqrt{\frac{1}{n}\sum_{i=1}^n R(\theta_i)^2}
\end{equation}
где $n$ — количество экспериментальных точек. Физический смысл RMSE — это «среднее» отклонение теоретической кривой от экспериментальных данных. Чем меньше значение RMSE, тем точнее модель описывает реальный физический процесс. Мы будем рассчитывать два значения: $\text{RMSE}_{\%}$ для линейного масштаба и $\text{RMSE}_{\text{dB}}$ для децибел.

\subsection{Интерактивная оценка и визуализация остатков}

В функции интерактивной подгонки мы реализовали сетку графиков $2 \times 2$ с помощью Matplotlib:
\begin{itemize}
    \item Слева вверху отображается линейная зависимость пропускания от угла поворота, а непосредственно под ней (слева внизу) — график остатков $R_{\%}(\theta)$.
    \item Справа вверху отображается логарифмическая зависимость (дБ), а под ней (справа внизу) — график остатков в дБ $R_{\text{dB}}(\theta)$.
\end{itemize}

\note{Графики остатков снабжены горизонтальной штриховой линией на уровне нуля. В идеальном случае, если бы наша модель абсолютно точно описывала эксперимент, все красные точки остатков лежали бы строго на этой нулевой линии. На практике они совершают колебания вокруг нуля из-за шумов измерения и неучтенных физических эффектов.}

Когда ты перемещаешь ползунки параметров ($P$, $D$, \texttt{period\_scale}, \texttt{loss\_factor}), программа в реальном времени:
\begin{enumerate}
    \item Пересчитывает теоретические спектры на экспериментальных углах.
    \item Находит новые массивы остатков $R_{\%}$ и $R_{\text{dB}}$ и мгновенно перерисовывает нижние графики.
    \item Вычисляет новые значения $\text{RMSE}_{\%}$ и $\text{RMSE}_{\text{dB}}$, выводя их прямо в заголовки верхних графиков.
\end{enumerate}
Это превращает подгонку в увлекательную игру: двигая ползунки, ты видишь, как колеблются точки остатков и как уменьшаются (или растут) цифры RMSE, стремясь к идеальному минимуму!

\subsection{Автоматизированная оптимизация параметров}

Хотя ручная интерактивная подгонка со слайдерами прекрасно развивает физическую интуицию, найти строгое математическое решение вручную крайне сложно — параметров четыре, и они взаимосвязаны. Чтобы автоматизировать этот процесс, мы применим методы математической оптимизации.

Наша цель — найти вектор параметров $\mathbf{x} = [P, D, \text{period\_scale}, \text{loss\_factor}]^T$, который минимизирует среднеквадратическую ошибку в децибелах:
\begin{equation}
\mathbf{x}^* = \arg\min_{\mathbf{x}} \text{RMSE}_{\text{dB}}(\mathbf{x})
\end{equation}

Почему для оптимизации мы выбираем именно логарифмическую шкалу (дБ)? Потому что значения пропускания по мощности охватывают несколько порядков (от $1.0$ до $10^{-4}$ и ниже). Если минимизировать линейную ошибку, алгоритм сосредоточится исключительно на идеальном совпадении максимумов пропускания, полностью проигнорировав область глубокой экстинкции, так как отклонения там численно ничтожны. Минимизация в шкале дБ обеспечивает одинаковое «внимание» алгоритма ко всем участкам кривой.

Для решения этой задачи мы воспользуемся функцией \texttt{scipy.optimize.minimize} и алгоритмом \textbf{Нелдера-Мида (Nelder-Mead, метод деформируемого многогранника)}.
\begin{itemize}
    \item \textbf{Почему Nelder-Mead?} Это прямой, безградиентный метод оптимизации. Он не требует вычисления производных, что делает его чрезвычайно устойчивым к шумам в данных и вычислительным погрешностям в сложных электродинамических формулах (таких как бесконечные ряды Blanco с пороговыми условиями).
\end{itemize}

В нашем скрипте \texttt{interactive\_fit.py} мы добавили кнопку \texttt{Button} с надписью «Оптимизировать». При её нажатии запускается следующий оптимизационный процесс:
\begin{minted}{python}
def loss_function(params):
    p, d, ps, ls = params
    if d >= p * ps:
        return 1e6  # Штраф за физически невозможную геометрию
    
    # Расчёт теории на экспериментальных углах
    y_db = np.array([theoretical.transmission_db_modified(
        ang, p, d, target_freq, ps, ls) for ang in exp_angles])
    
    # Возвращаем RMSE в дБ
    return np.sqrt(np.mean((exp_trans_db - y_db)**2))

res = minimize(loss_function, initial_guess, method='Nelder-Mead', bounds=bounds)
\end{minted}

Алгоритм использует текущие положения слайдеров в качестве начального приближения (\texttt{initial\_guess}), быстро спускается в многомерный минимум среднеквадратичной ошибки и автоматически сдвигает слайдеры в оптимальные положения.

\subsection{Анализ результатов подгонки}

Запустив скрипт и нажав кнопку «Оптимизировать», ты заметишь поразительную вещь. Ручная подгонка редко позволяет опустить $\text{RMSE}_{\text{dB}}$ ниже $2.5$ дБ. Алгоритм Нелдера-Мида находит глобальный минимум и снижает $\text{RMSE}_{\text{dB}}$ до значений менее $0.8$ дБ! При этом теоретическая кривая ложится на экспериментальные точки практически идеально во всем диапазоне углов.

Давай проанализируем физический смысл найденных оптимальных параметров:
\begin{itemize}
    \item \textbf{Шаг решетки $P$ и масштабный множитель \texttt{period\_scale}:} Паспортное значение шага составляет $12.5$ мкм. Оптимизатор обычно находит \texttt{period\_scale} $\approx 1.015 - 1.03$, что соответствует эффективному шагу $12.7 - 12.9$ мкм. Это абсолютно логично с физической точки зрения — при намотке микропроволоки на рамку поляризатора она испытывает сильное механическое натяжение, что приводит к упругому растяжению и незначительному увеличению шага решетки.
    
    \item \textbf{Диаметр проволоки $D$:} Оптимизированный диаметр проволоки обычно близок к паспортным $5.0$ мкм (в пределах погрешности микрометра).
    
    \item \textbf{Коэффициент потерь \texttt{loss\_factor}:} Оптимизатор находит ненулевое значение \texttt{loss\_factor} $\approx 0.1 - 0.25$ дБ/ТГц. Без этого параметра теоретическая кривая в области скрещенных поляризаторов уходила бы в бесконечный ноль (минус бесконечность в дБ), тогда как в реальности всегда присутствует фоновое рассеяние терагерцовых волн на микронеровностях проволоки и омические потери в металле. Ненулевой \texttt{loss\_factor} позволяет идеально сгладить и прижать «дно» теоретической кривой к уровню реальной экстинкции эксперимента.
\end{itemize}

\begin{minitaskbox}
\textbf{Мини-исследование: Влияние начального приближения}
\\
Попробуй сдвинуть слайдеры в заведомо плохие положения (например, сделай $P = 30$ мкм, $D = 1$ мкм) и нажми кнопку «Оптимизировать». Справляется ли алгоритм Нелдера-Мида с поиском глобального минимума из плохого стартового приближения? Запиши полученные значения параметров и сравни их со значениями, полученными при хорошем старте. Подумай, почему в сложных задачах важно начинать с физически обоснованного начального приближения.
\end{minitaskbox}

\subsection{Обновление файла config.py}
Хотя в коде выше мы использовали начальные значения напрямую, правильным тоном будет добавить их в наш файл \texttt{config.py}, чтобы они стали доступны другим скриптам в будущем.
Добавь следующие параметры в свой \mycode{config.py}: \mycode{P_DEFAULT = 12.5e-6} (паспортный шаг решетки, м), \mycode{D_DEFAULT = 5.0e-6} (паспортный диаметр проволоки, м), \mycode{PERIOD_SCALE_DEFAULT = 1.0} (масштабный множитель периода) и \mycode{LOSS_FACTOR_DEFAULT = 0.0} (коэффициент потерь).

\subsection{Итоги раздела}
Мы перешли от субъективного визуального сопоставления теории и эксперимента к строгим методам математической оптимизации. Введение остатков в процентах и децибелах позволило детально локализовать погрешности модели на разных масштабах пропускания, а метрика RMSE дала нам числовой критерий качества. Наконец, интеграция безградиентного метода Нелдера-Мида прямо в интерактивное окно позволила автоматически и с прецизионной точностью находить реальные физические параметры поляризаторов. В следующем разделе мы масштабируем этот подход и проведём автоматизированный анализ параметров решётки по всему частотному спектру!
